\section{Dimension}
\label{section-dimension}

\noindent
Please compare with
Topology, Section \ref{topology-section-krull-dimension}.

\begin{definition}
\label{definition-chain-primes}
Let $R$ be a ring. A {\it chain of prime ideals} is a sequence
$\mathfrak p_0 \subset \mathfrak p_1 \subset \ldots \subset \mathfrak p_n$
of prime ideals of $R$ such that $\mathfrak p_i \not = \mathfrak p_{i + 1}$
for $i = 0, \ldots, n - 1$. The {\it length} of this chain of prime
ideals is $n$.
\end{definition}

\noindent
Recall that we have an inclusion reversing bijection between prime
ideals of a ring $R$ and irreducible closed subsets of $\Spec(R)$,
see Lemma \ref{lemma-irreducible}.

\begin{definition}
\label{definition-Krull}
The {\it Krull dimension} of the ring $R$ is the
Krull dimension of the topological space $\Spec(R)$, see
Topology, Definition \ref{topology-definition-Krull}.
In other words it is the supremum of the integers $n\geq 0$
such that $R$ has a chain of prime ideals
$$
\mathfrak p_0
\subset
\mathfrak p_1
\subset
\ldots
\subset
\mathfrak p_n, \quad
\mathfrak p_i \not = \mathfrak p_{i + 1}.
$$
of length $n$.
\end{definition}

\begin{definition}
\label{definition-height}
The {\it height} of a prime ideal $\mathfrak p$ of
a ring $R$ is the dimension of the local ring $R_{\mathfrak p}$.
\end{definition}

\begin{lemma}
\label{lemma-dimension-height}
The Krull dimension of $R$ is the supremum of the
heights of its (maximal) primes.
\end{lemma}

\begin{proof}
This is so because we can always add a maximal ideal at the end of a chain
of prime ideals.
\end{proof}

\begin{lemma}
\label{lemma-Noetherian-dimension-0}
A Noetherian ring of dimension $0$ is Artinian.
Conversely, any Artinian ring is Noetherian of dimension zero.
\end{lemma}

\begin{proof}
Assume $R$ is a Noetherian ring of dimension $0$.
By Lemma \ref{lemma-Noetherian-topology} the space $\Spec(R)$
is Noetherian. By Topology, Lemma \ref{topology-lemma-Noetherian} we see
that $\Spec(R)$ has finitely many irreducible
components, say $\Spec(R) = Z_1 \cup \ldots \cup Z_r$.
According to Lemma \ref{lemma-irreducible} each $Z_i = V(\mathfrak p_i)$
with $\mathfrak p_i$ a minimal ideal. Since the dimension is $0$
these $\mathfrak p_i$ are also maximal. Thus $\Spec(R)$
is the discrete topological space with elements $\mathfrak p_i$.
All elements $f$ of the Jacobson radical $\bigcap \mathfrak p_i$
are nilpotent since otherwise $R_f$ would not be the zero ring
and we would have another prime.
By Lemma \ref{lemma-product-local} $R$ is equal to
$\prod R_{\mathfrak p_i}$. Since $R_{\mathfrak p_i}$
is also Noetherian and dimension $0$, the previous arguments
show that its radical $\mathfrak p_iR_{\mathfrak p_i}$ is locally nilpotent.
Lemma \ref{lemma-Noetherian-power} gives
$\mathfrak p_i^nR_{\mathfrak p_i} = 0$ for some $n \geq 1$.
By Lemma \ref{lemma-length-finite} we conclude that $R_{\mathfrak p_i}$
has finite length over $R$. Hence we conclude that $R$
is Artinian by Lemma \ref{lemma-artinian-finite-length}.

\medskip\noindent
If $R$ is an Artinian ring then by Lemma \ref{lemma-artinian-finite-length}
it is Noetherian. All of its primes are maximal by a combination
of Lemmas \ref{lemma-artinian-finite-nr-max},
\ref{lemma-artinian-radical-nilpotent} and \ref{lemma-product-local}.
\end{proof}

\noindent
In the following we will use the invariant $d(-)$ defined
in Definition \ref{definition-d}. Here is a warm up lemma.

\begin{lemma}
\label{lemma-dimension-0-d-0}
Let $R$ be a Noetherian local ring.
Then $\dim(R) = 0 \Leftrightarrow d(R) = 0$.
\end{lemma}

\begin{proof}
This is because $d(R) = 0$ if and only if $R$ has finite
length as an $R$-module. See Lemma \ref{lemma-artinian-finite-length}.
\end{proof}

\begin{proposition}
\label{proposition-dimension-zero-ring}
Let $R$ be a ring. The following are equivalent:
\begin{enumerate}
\item $R$ is Artinian,
\item $R$ is Noetherian and $\dim(R) = 0$,
\item $R$ has finite length as a module over itself,
\item $R$ is a finite product of Artinian local rings,
\item $R$ is Noetherian and $\Spec(R)$ is a
finite discrete topological space,
\item $R$ is a finite product of Noetherian local rings
of dimension $0$,
\item $R$ is a finite product of Noetherian local rings
$R_i$ with $d(R_i) = 0$,
\item $R$ is a finite product of Noetherian local rings
$R_i$ whose maximal ideals are nilpotent,
\item $R$ is Noetherian, has finitely many maximal
ideals and its Jacobson radical ideal is nilpotent, and
\item $R$ is Noetherian and there are no strict inclusions
among its primes.
\end{enumerate}
\end{proposition}

\begin{proof}
This is a combination of Lemmas
\ref{lemma-product-local},
\ref{lemma-artinian-finite-length},
\ref{lemma-Noetherian-dimension-0}, and
\ref{lemma-dimension-0-d-0}.
\end{proof}

\begin{lemma}
\label{lemma-height-1}
Let $R$ be a local Noetherian ring.
The following are equivalent:
\begin{enumerate}
\item
\label{item-dim-1}
$\dim(R) = 1$,
\item
\label{item-d-1}
$d(R) = 1$,
\item
\label{item-Vx}
there exists an $x \in \mathfrak m$, $x$ not nilpotent
such that $V(x) = \{\mathfrak m\}$,
\item
\label{item-x}
there exists an $x \in \mathfrak m$, $x$ not nilpotent
such that $\mathfrak m = \sqrt{(x)}$, and
\item
\label{item-ideal-1}
there exists an ideal of definition generated by $1$ element,
and no ideal of definition is generated by $0$ elements.
\end{enumerate}
\end{lemma}

\begin{proof}
First, assume that $\dim(R) = 1$.
Let $\mathfrak p_i$ be the minimal primes of $R$.
Because the dimension is $1$ the only other prime of $R$
is $\mathfrak m$.
According to Lemma \ref{lemma-Noetherian-irreducible-components}
there are finitely many. Hence we can find $x \in \mathfrak m$,
$x \not \in \mathfrak p_i$, see Lemma \ref{lemma-silly}.
Thus the only prime containing $x$ is $\mathfrak m$ and
hence (\ref{item-Vx}).

\medskip\noindent
If (\ref{item-Vx}) then $\mathfrak m = \sqrt{(x)}$ by
Lemma \ref{lemma-Zariski-topology}, and hence (\ref{item-x}).
The converse is clear as well.
The equivalence of (\ref{item-x}) and (\ref{item-ideal-1}) follows
from directly the definitions.

\medskip\noindent
Assume (\ref{item-ideal-1}).
Let $I = (x)$ be an ideal of definition.
Note that $I^n/I^{n + 1}$ is a quotient of $R/I$ via multiplication
by $x^n$ and hence $\text{length}_R(I^n/I^{n + 1})$ is bounded.
Thus $d(R) = 0$ or $d(R) = 1$, but $d(R) = 0$ is excluded
by the assumption that $0$ is not an ideal of definition.

\medskip\noindent
Assume (\ref{item-d-1}). To get a contradiction, assume there
exist primes $\mathfrak p \subset \mathfrak q \subset \mathfrak m$,
with both inclusions strict. Pick some ideal of definition $I \subset R$.
We will repeatedly use
Lemma \ref{lemma-hilbert-ses-chi}. First of all
it implies, via the exact sequence
$0 \to \mathfrak p \to R \to R/\mathfrak p \to 0$,
that $d(R/\mathfrak p) \leq 1$. But it clearly cannot
be zero. Pick $x\in \mathfrak q$, $x\not \in \mathfrak p$.
Consider the short exact sequence
$$
0 \to R/\mathfrak p \to R/\mathfrak p \to R/(xR + \mathfrak p) \to 0.
$$
This implies that $\chi_{I, R/\mathfrak p} - \chi_{I, R/\mathfrak p}
- \chi_{I, R/(xR + \mathfrak p)} = - \chi_{I, R/(xR + \mathfrak p)}$
has degree $ < 1$. In other words, $d(R/(xR + \mathfrak p)) = 0$,
and hence $\dim(R/(xR + \mathfrak p)) = 0$, by
Lemma \ref{lemma-dimension-0-d-0}. But $R/(xR + \mathfrak p)$
has the distinct primes $\mathfrak q/(xR + \mathfrak p)$ and
$\mathfrak m/(xR + \mathfrak p)$ which gives the desired contradiction.
\end{proof}

\begin{proposition}
\label{proposition-dimension}
Let $R$ be a local Noetherian ring. Let $d \geq 0$ be an integer.
The following are equivalent:
\begin{enumerate}
\item
\label{item-dim-d}
$\dim(R) = d$,
\item
\label{item-d-d}
$d(R) = d$,
\item
\label{item-ideal-d}
there exists an ideal of definition generated by $d$ elements,
and no ideal of definition is generated by fewer than $d$ elements.
\end{enumerate}
\end{proposition}

\begin{proof}
This proof is really just the same as the proof of Lemma
\ref{lemma-height-1}. We will prove the proposition by induction
on $d$. By Lemmas \ref{lemma-dimension-0-d-0} and \ref{lemma-height-1}
we may assume that $d > 1$. Denote the minimal number of
generators for an ideal of definition of $R$ by $d'(R)$.
We will prove the inequalities
$\dim(R) \geq d'(R) \geq d(R) \geq \dim(R)$,
and hence they are all equal.

\medskip\noindent
First, assume that $\dim(R) = d$.
Let $\mathfrak p_i$ be the minimal primes of $R$.
According to Lemma \ref{lemma-Noetherian-irreducible-components}
there are finitely many. Hence we can find $x \in \mathfrak m$,
$x \not \in \mathfrak p_i$, see Lemma \ref{lemma-silly}.
Note that every maximal chain of primes starts with some $\mathfrak p_i$,
hence the dimension of $R/xR$ is at most $d-1$. By induction
there are $x_2, \ldots, x_d$ which generate an ideal of definition
in $R/xR$. Hence $R$ has an ideal of definition generated
by (at most) $d$ elements.

\medskip\noindent
Assume $d'(R) = d$. Let $I = (x_1, \ldots, x_d)$ be an ideal
of definition. Note that $I^n/I^{n + 1}$ is a quotient of a direct
sum of $\binom{d + n - 1}{d - 1}$ copies $R/I$ via multiplication
by all degree $n$ monomials in $x_1, \ldots, x_d$.
Hence $\text{length}_R(I^n/I^{n + 1})$ is bounded by a polynomial
of degree $d-1$. Thus $d(R) \leq d$.

\medskip\noindent
Assume $d(R) = d$. Consider a chain of primes
$\mathfrak p \subset \mathfrak q \subset
\mathfrak q_2 \subset \ldots \subset \mathfrak q_e = \mathfrak m$,
with all inclusions strict, and $e \geq 2$.
Pick some ideal of definition $I \subset R$.
We will repeatedly use
Lemma \ref{lemma-hilbert-ses-chi}. First of all
it implies, via the exact sequence
$0 \to \mathfrak p \to R \to R/\mathfrak p \to 0$,
that $d(R/\mathfrak p) \leq d$. But it clearly cannot
be zero. Pick $x\in \mathfrak q$, $x\not \in \mathfrak p$.
Consider the short exact sequence
$$
0 \to R/\mathfrak p \to R/\mathfrak p \to R/(xR + \mathfrak p) \to 0.
$$
This implies that $\chi_{I, R/\mathfrak p} - \chi_{I, R/\mathfrak p}
- \chi_{I, R/(xR + \mathfrak p)} = - \chi_{I, R/(xR + \mathfrak p)}$
has degree $ < d$. In other words, $d(R/(xR + \mathfrak p)) \leq d - 1$,
and hence $\dim(R/(xR + \mathfrak p)) \leq d - 1$, by
induction. Now $R/(xR + \mathfrak p)$ has the chain of prime ideals
$\mathfrak q/(xR + \mathfrak p) \subset \mathfrak q_2/(xR + \mathfrak p)
\subset \ldots \subset \mathfrak q_e/(xR + \mathfrak p)$ which gives
$e - 1 \leq d - 1$. Since we started with an arbitrary chain of
primes this proves that $\dim(R) \leq d(R)$.

\medskip\noindent
Reading back the reader will see we proved the circular
inequalities as desired.
\end{proof}

\noindent
Let $(R, \mathfrak m)$ be a Noetherian local ring.
From the above it is clear that $\mathfrak m$ cannot be
generated by fewer than $\dim(R)$ variables.
By Nakayama's Lemma \ref{lemma-NAK} the minimal number
of generators of $\mathfrak m$ equals $\dim_{\kappa(\mathfrak m)}
\mathfrak m/\mathfrak m^2$. Hence we have the following
fundamental inequality
$$
\dim(R) \leq \dim_{\kappa(\mathfrak m)} \mathfrak m/\mathfrak m^2.
$$
It turns out that the rings where equality holds
have a lot of good properties. They are called
regular local rings.

\begin{definition}
\label{definition-regular-local}
Let $(R, \mathfrak m)$ be a Noetherian local ring of dimension $d$.
\begin{enumerate}
\item A {\it system of parameters of $R$} is a sequence of elements
$x_1, \ldots, x_d \in \mathfrak m$ which generates an ideal of
definition of $R$,
\item if there exist $x_1, \ldots, x_d \in \mathfrak m$
such that $\mathfrak m = (x_1, \ldots, x_d)$ then we call
$R$ a {\it regular local ring} and $x_1, \ldots, x_d$ a {\it regular
system of parameters}.
\end{enumerate}
\end{definition}

\noindent
The following lemmas are clear from the proofs of the
lemmas and proposition above, but we spell them out so we have
convenient references.

\begin{lemma}
\label{lemma-minimal-over-1}
Let $R$ be a Noetherian ring. Let $x \in R$.
\begin{enumerate}
\item If $\mathfrak p$ is minimal over $(x)$
then the height of $\mathfrak p$ is $0$ or $1$.
\item If $\mathfrak p, \mathfrak q \in \Spec(R)$ and $\mathfrak q$
is minimal over $(\mathfrak p, x)$, then there is no prime strictly
between $\mathfrak p$ and $\mathfrak q$.
\end{enumerate}
\end{lemma}

\begin{proof}
Proof of (1). If $\mathfrak p$ is minimal over $x$, then the only
prime ideal of $R_\mathfrak p$ containing $x$ is the maximal ideal
$\mathfrak p R_\mathfrak p$. This is true because the primes of
$R_\mathfrak p$ correspond $1$-to-$1$ with the primes of $R$ contained
in $\mathfrak p$, see Lemma \ref{lemma-spec-localization}.
Hence Lemma \ref{lemma-height-1} shows $\dim(R_\mathfrak p) = 1$
if $x$ is not nilpotent in $R_\mathfrak p$. Of course, if
$x$ is nilpotent in $R_\mathfrak p$ the argument gives that
$\mathfrak pR_\mathfrak p$ is the only prime ideal and we see
that the height is $0$.

\medskip\noindent
Proof of (2). By part (1) we see that $\mathfrak q/\mathfrak p$
is a prime of height $1$ or $0$ in $R/\mathfrak p$. This immediately
implies there cannot be a prime strictly between $\mathfrak p$
and $\mathfrak q$.
\end{proof}

\begin{lemma}
\label{lemma-minimal-over-r}
Let $R$ be a Noetherian ring. Let $f_1, \ldots, f_r \in R$.
\begin{enumerate}
\item If $\mathfrak p$ is minimal over $(f_1, \ldots, f_r)$
then the height of $\mathfrak p$ is $\leq r$.
\item If $\mathfrak p, \mathfrak q \in \Spec(R)$ and
$\mathfrak q$ is minimal over $(\mathfrak p, f_1, \ldots, f_r)$,
then every chain of primes between $\mathfrak p$ and $\mathfrak q$
has length at most $r$.
\end{enumerate}
\end{lemma}

\begin{proof}
Proof of (1). If $\mathfrak p$ is minimal over $f_1, \ldots, f_r$,
then the only prime ideal of $R_\mathfrak p$ containing $f_1, \ldots, f_r$
is the maximal ideal $\mathfrak p R_\mathfrak p$. This is true because
the primes of $R_\mathfrak p$ correspond $1$-to-$1$ with the primes of
$R$ contained in $\mathfrak p$, see Lemma \ref{lemma-spec-localization}.
Hence Proposition \ref{proposition-dimension} shows
$\dim(R_\mathfrak p) \leq r$.

\medskip\noindent
Proof of (2). By part (1) we see that $\mathfrak q/\mathfrak p$
is a prime of height $\leq r$. This immediately
implies the statement about chains of primes between $\mathfrak p$
and $\mathfrak q$.
\end{proof}

\begin{lemma}
\label{lemma-one-equation}
Suppose that $R$ is a Noetherian local ring and $x\in \mathfrak m$ an
element of its maximal ideal. Then $\dim R \leq \dim R/xR + 1$.
If $x$ is not contained in any of the minimal primes of $R$
then equality holds. (For example if $x$ is a nonzerodivisor.)
\end{lemma}

\begin{proof}
If $x_1, \ldots, x_{\dim R/xR} \in R$ map to elements of $R/xR$ which
generate an ideal of definition for $R/xR$, then $x, x_1, \ldots,
x_{\dim R/xR}$ generate an ideal of definition for $R$. Hence
the inequality by Proposition \ref{proposition-dimension}.
On the other hand, if $x$ is not contained in any minimal
prime of $R$, then the chains of primes in $R/xR$ all give
rise to chains in $R$ which are at least one step away
from being maximal.
\end{proof}

\begin{lemma}
\label{lemma-elements-generate-ideal-definition}
Let $(R, \mathfrak m)$ be a Noetherian local ring.
Suppose $x_1, \ldots, x_d \in \mathfrak m$ generate an
ideal of definition and $d = \dim(R)$. Then
$\dim(R/(x_1, \ldots, x_i)) = d - i$ for all $i = 1, \ldots, d$.
\end{lemma}

\begin{proof}
Follows either from the proof of Proposition \ref{proposition-dimension},
or by using induction on $d$ and Lemma \ref{lemma-one-equation}.
\end{proof}







